\section{Introduction}
\label{sec:introduction}

Bounded representations of zero are a standard finite-state object in
Pisot numeration.  Their regularity supports normalization and arithmetic
algorithms, but regularity by itself does not control how long an
aperture-dependent chain of overlapping zero representations can remain
transient.  We consider the following narrow quantitative question.  Fix a
Pisot recurrence $U=(u_j)$ and a coefficient alphabet $[-D,D]\cap\mathbb Z$.
At aperture $m$, use a length-$m$ coefficient word as a state transition
exactly when it can be completed by a coefficient of the distinguished
weight $u_m$ to represent zero.  Consecutive edges overlap by one position.
How long can a path beginning with an exposed nonzero coefficient survive if
it never reaches a cycle?

Our first main result, Theorem~\ref{thm:bounded-zero-linear-transient}, gives
an effective linear answer.  There is a computable $C(U,D)<\infty$ such that
every such transient has fewer than $C(U,D)m$ edges.  Equivalently, after a
nonzero chain is anchored at its first exposed nonzero coefficient, it
either enters a repeatable periodic tail or has only linearly many further
overlaps.  The coefficient completing an edge is not assumed to belong to
$[-D,D]$; it is proved to lie in a fixed interval
$[-K_{U,D},K_{U,D}]$.  This mixed-alphabet point is essential to the correct
edge formulation.

The mechanism is local.  Lemma
\ref{lem:bounded-zero-arbitrary-alphabet-collapse} compares two consecutive
zero representations under every algebraic embedding.  The contracting
conjugates bound the nondominant sums, while a finite exact separation in the
dominant embedding prevents any nonzero value of
$e_{t+m}+k_t-\beta k_{t+1}$ once $m$ is large enough.  Irrationality then
forces
\[
k_{t+1}=0,\qquad e_{t+m}=-k_t.
\]
Along a path, every quotient after the first vanishes and all later appended
coefficients vanish after one exceptional position.  A path of $m+1$ edges
therefore reaches the zero state, which already has a loop.  The finitely
many smaller apertures are handled by exact acyclic longest-path searches.
This gives an effective constant without counting all
$(2D+1)^{m-1}$ states.

The proof requires neither Condition~F nor preservation of leading zeros.
Finite-state recognition of bounded zero representations in Pisot systems is
classical \cite{Frougny1992}, and recent work places zero preservation in a
precise Condition~F framework \cite{CartonSudberyYassawi2026}.  Those results
do not imply an aperture-linear survival bound for the moving overlap graphs
considered here.  Conversely, our result does not assert finite expansion,
normalization complexity, or a general synchronization theorem.  The
distinguished terminal weight $u_m$ and consecutive recurrence weights are
part of the statement.

The second part applies the standard-object theorem to a symbolic recoding.
For a strictly increasing Pisot numeration system, order the legal
length-$m$ words by canonical rank, extend the same positional rank to the
full digit cube, reduce it modulo the number $u_m$ of legal words, and return
the legal word with that residue.  Overlap these recoded windows along the
shift.  We call the resulting factor code a cyclic language-rank recoding.
It is a congruence of integer language ranks, not numerical
$\beta$-normalization in $\mathbb Z[\beta]$.

Lemma~\ref{lem:pisot-exact-collision-quotient} proves that the
coordinatewise-difference quotient of the pair graph loses no collision:
every bounded digit difference is realized by its positive and negative
parts.  The quotient is exactly the graph $G_m(U,d_U)$ from the first part.
The longest path from a first-coordinate-nonzero state is exactly one less
than the least future-only inverse window, and injectivity is equivalent to
the absence of a reachable cycle.  Corollary
\ref{cor:pisot-linear-causal-depth} therefore gives
\[
\Phi_{U,m}\text{ injective}
\quad\Longrightarrow\quad
\ell_{\rm cau}(U,m)\le C_Um. \tag{A1}
\]
This is a consequence of the bounded-zero theorem and the exact quotient,
not the definitional center of the argument.

Classical pair and fiber graphs supply reachable-cycle tests and finite
inverse windows for injective local maps
\cite{AmorosoPatt1972,Richardson1972,Nasu1977,Head1989,LindMarcus1995,
Nasu1995}.  Their generic state-count bound is exponential in $m$ here.
Likewise, Ashley's decoder-window estimates are linear in the number of
states of a presenting graph, not in this aperture
\cite{Ashley1988,Ashley1996}.  The bounded-delay theorem of
Frougny and Sakarovitch starts from a rational relation already known to have
bounded delay; it does not derive an $O(m)$ delay from injectivity
\cite{FrougnySakarovitch1991}.  The arithmetic adjacent-collapse lemma is
what changes the scale from exponential to linear.

Finally, Theorem~\ref{thm:fixed-cubic-sharpness} proves sharpness inside one
fixed system.  Let $\theta$ be the leading root of
$x^3-2x^2+x-1$ and use the binary recurrence beginning $1,2,4$.  For every
$m\ge4$,
\[
\Phi_{\theta,m}\text{ is injective},\qquad
\ell_{\rm cau}(\theta,m)=2\lfloor m/2\rfloor-1. \tag{A2}
\]
The proof classifies the penultimate obstruction set exactly as
$\{E_m,-E_m\}$ and the next obstruction set as empty.  Read only as bounded
coefficient words, the same $E_m$ give paths of length at least $m-3$ in
$G_m(U_\theta,1)$.  One object therefore proves sharpness both for the
inverse-depth corollary and for the bounded-zero transient theorem.

The article is organized in this logical order.  Section
\ref{sec:bounded-zero-transients} defines the overlap chains, proves the
arbitrary-alphabet collapse lemma, and proves the linear transient theorem.
Section~\ref{sec:cyclic-rank-application} introduces cyclic ranks, proves the
exact collision quotient, and deduces \textup{(A1)}.  Section
\ref{sec:fixed-cubic-sharpness} proves \textup{(A2)}--\textup{(A3)} and the
terminal classification.  Section~\ref{sec:comparative-context} records the
limited comparison with the family-classification problem.  The companion
article by the present authors, \emph{Exact overlap thresholds and
future-only inverse depth for cyclic ranks of quadratic and simple-Parry
beta-languages}, contains the full Fibonacci, quadratic, and simple-Parry
classification.
